\section{A paired controlled-continuum construction}\label{sec:continuum48}
Consider the same Borel problem \eqref{eq:target48}, with transition kernels
\begin{equation}\label{eq:atoms48}
 P_t^a(x,\cdot)=\sum_{s=1}^{S}\omega_{tas}(x)\delta_{F_{tas}(x)}(\cdot),
 \qquad \omega_{tas}\ge0,\quad\sum_s\omega_{tas}=1,
\end{equation}
where the maps and weights are Borel. They may be action dependent, noninvertible, and state dependent. Let the \emph{objective} initial law have finite support $H_0$, and recursively set
\begin{equation}\label{eq:closure48}
 H_{t+1}=\bigcup_{x\in H_t,a,s:\omega_{tas}(x)>0}\{F_{tas}(x)\}.
\end{equation}
Equal successors are identified, not copied into history-specific nodes. This is essential: two paths that reach the same state and date must use the same future policy.

\begin{theorem}[Finite observation and global extension]\label{thm:closure48}
Under Section~\ref{sec:model48} and \eqref{eq:atoms48}, with a finitely supported objective initial law, the unrestricted Borel common-policy value equals the finite common-policy value on the forward sets \eqref{eq:closure48}:
\begin{equation}\label{eq:exactclosure48}
 \KR^X(\nu,\varepsilon)=\KR^H(\nu,\varepsilon).
\end{equation}
Every feasible finite policy has an explicit Borel extension satisfying the operating constraint at \emph{every} point of $X$, not just at the forward sets. Therefore a paired finite certificate $[L,U]$ is a paired certificate for the full-state problem with exactly the same endpoints. The construction uses at most $|H_0|(mS)^t$ states at date $t$ and incurs zero spatial approximation error.
\end{theorem}
\begin{proof}
All possible successors of a point of $H_t$ belong to $H_{t+1}$, so the operating Bellman value on these sets equals the restriction of $V^X$ by backward induction. Every globally feasible policy restricts to a feasible finite policy and has the same objective. This proves $\KR^H\le\KR^X$.

For the reverse direction take any finite feasible policy $p^H$. Put $\widehat J_T=g$. After defining the later extension, set
\begin{equation}\label{eq:extension48}
 \widehat p_t(\cdot\mid x)=
 \begin{cases}
 p_t^H(\cdot\mid x),&x\in H_t,\\
 \delta_{a_t(x)},&x\notin H_t,
 \end{cases}\qquad
 a_t(x)=\text{first maximizer of }r_t(x,a)+\beta P_t^a\widehat J_{t+1}(x).
\end{equation}
Finite sets are Borel. Finite-action first-maximizer selection and the Borel kernels make the extension Borel. Values on $H_t$ are exactly the finite values because all successors stay on the forward sets. They satisfy $\widehat J_t\ge V_t-\varepsilon$. Off $H_t$, assuming global feasibility at date $t+1$, one has
\[
 \widehat J_t(x)=\max_a\{r_t(x,a)+\beta P_t^a\widehat J_{t+1}(x)\}
 \ge V_t(x)-\beta\varepsilon\ge V_t(x)-\varepsilon.
\]
This completes the induction. Starting under $\nu$, the policy never leaves the finite forward sets, so its implementation objective is unchanged. Taking infima proves the reverse inequality. Each state has at most $mS$ successors, which gives the size bound. No state or probability is perturbed, establishing the zero spatial error.
\end{proof}

\begin{corollary}[Executable paired accuracy]\label{cor:continuum48}
Suppose the finite forward model can be constructed with exact rational arithmetic, including decidable equality of successor states. Apply Theorem~\ref{thm:finite48} to that model and use \eqref{eq:extension48} off the forward sets. This gives a computable paired hierarchy and a stopping rule $U-L\le\eta$ for \eqref{eq:target48}. A running maximum of certified lower endpoints and running minimum of feasible upper endpoints give monotone sequences converging to the same value when the work caps are removed and targets tend to zero.
\end{corollary}

The extension is also operational for the affine rational model below. At any queried state, evaluate its finitely branching continuation and test membership in the finite rational forward sets. This computes an action without an exact global operating oracle. In a piecewise-affine implementation, ordinary max-affine breakpoints coexist with exceptional preimages of the assigned finite policies. Both counts may grow exponentially with the horizon. The theorem guarantees a representation and a matched certificate, not cheap symbolic compilation. The publication contract checker independently reconstructs all forward states, merges, transition rows, rewards, and the terminal folding used in the experiment.

\subsection{An observed-fleet service contract with original controlled atoms}
The continuous condition is $x\in[0,1]$. Operating rewards are $r(x,a)=x-c_a$, with
\[
 (c_0,c_1,c_2)=(0,3/20,9/20),\quad g(x)=x/2,\quad\beta=19/20.
\]
An installed action $0$ costs nothing to retain; either departure costs $1+x$. The two shock outcomes for the three actions are
\begin{align*}
 F_{0}(x)&\in\{3x/4,\;3x/4+1/20\},\\
 F_{1}(x)&\in\{17x/20+1/10,\;17x/20+7/50\},\\
 F_{2}(x)&\in\{x/10+4/5,\;x/10+17/20\},
\end{align*}
with probabilities $3/5,2/5$. These are the original affine controlled atoms, not a finite-state replacement for different dynamics. Condition affects reward, revision expense, and every controlled successor. We fix $\varepsilon=1/100$ and the observed initial fleet law $\nu=(3/4)\delta_{1/5}+(1/4)\delta_{4/5}$, at horizons 2 and 3.

A service organization can interpret the statewise allowance as a warranty on remaining operating value whenever its rule is restarted, including off-fleet states; the objective weights the vehicles in the observed fleet. This is a specified decision contract, not evidence that any particular organization uses it. The extension in Theorem~\ref{thm:closure48} is what makes the warranty global despite the finite objective population. Lottery administration is assumed permitted by that contract and is charged separately if nonzero.

The two certified optimal costs are approximately 1.2696 and 2.8687 normalized implementation units. Hence a two-unit budget permits the two-period task but not the three-period task under the stated operating guarantee. This is an internally verified budget distinction in designed primitives, not an estimate of market maintenance costs or empirical welfare. Multiplying the cost scale and budget by the same positive constant preserves the conclusion. The earlier arbitrary monetary conversion is retained only in the clearly identified historical manuscript.

For horizon 2 the active forward-set sizes are $2,11,63$; for horizon 3 they are $2,11,63,367$. The last transition's expected $g$ is folded into the last-period reward exactly, avoiding terminal padding. The policy problem has respectively 13 and 76 active decision nodes. Dummy nodes used for rectangular storage carry zero costs and have no incoming transitions from active states. Their equations are checked separately; they do not stand for omitted physical states. Table~\ref{tab:fleet48} reports the paired intervals.

\subsection{What is not an initial-law approximation}
Equation~\eqref{eq:exactclosure48} keeps $\nu$ fixed. It does not authorize replacing a uniform initial law by finitely many atoms. In particular, the earlier thirty positive-cost uniform-initial-law intervals are unchanged, and the earlier nonlinear two-state widths remain roughly 25--76\% at allowance $1/2$. Their original results and proofs are reproduced in the retained technical material. The one-sided Borel density theorem likewise remains a one-sided theorem outside the hypotheses of Theorem~\ref{thm:closure48}. A general matched numerical hierarchy for diffuse initial laws with controlled moving atoms requires further structure and is not inferred from finite observation.
